\section*{अध्याय \ref{ch:b1:earthquake-source} --- भूकंप का उद्गम}

\begin{solution}{exo:b1:earthquake-source:1}
(a) $\log_{10}(\num{4e17}) = 17.60$, इसलिए $M_w = \tfrac23(17.60 - 9.1) = 5.67$।
(b) $\log_{10}(\num{5.59e22}) = 22.75$ और $M_w = \tfrac23(22.75 - 9.1) =
9.10$। (c) $M_0 = 10^{1.5\times7.0 + 9.1} = 10^{19.6} = \qty{3.98e19}{N.m}$।
\end{solution}

\begin{solution}{exo:b1:earthquake-source:2}
(a) $T = D/\dot s = \qty{3}{m}/\qty{0.025}{m.yr^{-1}} = \qty{120}{yr}$।
(b) $c = 0.5$ के साथ, $T = D/(c\dot s) = \qty{240}{yr}$: विसर्पण आधी गति लेता है, और उसी सरकन के भूकंप आधी बार आते हैं।
\end{solution}

\begin{solution}{exo:b1:earthquake-source:3}
\cref{prop:b1:earthquake-source:moment-steps} से आघूर्णों का अनुपात
$10^{1.5\times(7.4 - 6.2)} = 10^{1.8} = 63.1$ है: परिमाण 7.4 के एक भूकंप का आघूर्ण मुक्त करने के लिए परिमाण 6.2 के 63 भूकंप लगते हैं।
\end{solution}

\begin{solution}{exo:b1:earthquake-source:4}
$b = 1$ के साथ, परिमाण की हर इकाई गिनती को 10 से भाग देती है: परिमाण 5 या उससे अधिक के प्रति वर्ष 10 भूकंप, परिमाण 6 या उससे अधिक के 1, और परिमाण 7 या उससे अधिक के 0.1, यानी औसतन हर 10 वर्ष में एक।
\end{solution}

\begin{solution}{exo:b1:earthquake-source:5}
(a) $A = \num{4.3e5} \times \qty{1.5e4}{m} = \qty{6.45e9}{m^2}$, और
$M_0 = \num{3e10}\times\num{6.45e9}\times4 = \qty{7.74e20}{N.m}$;
$\log_{10}M_0 = 20.89$ और $M_w = \tfrac23(20.89 - 9.1) = 7.86$। (b) \qty{10}{km} की गहराई क्षेत्रफल और आघूर्ण को 1.5 से विभाजित करती है, जो परिमाण को
$\tfrac23\log_{10}1.5 = 0.12$ से घटाती है, 7.74 तक। (c) पहला, 7.86, सूची के 7.9 तक पूर्णांकित होता है; दूसरा थोड़ा कम है। परिमाण गहराई के प्रति बहुत संवेदनशील नहीं है: क्षेत्रफल पर 1.5 का गुणक केवल 0.12 परिमाण है।
\end{solution}

\begin{solution}{exo:b1:earthquake-source:6}
(a) $D = M_0/(\mu A) = \num{5.59e22}/(\num{3e10}\times\num{4e5}\times\num{1.5e5})
= \qty{31.1}{m}$। (b) $\qty{31.1}{m}/\qty{0.083}{m.yr^{-1}} = \qty{374}{yr}$।
(c) अवरुद्ध भाग 0.8 हो, तो $31.1/(0.8\times0.083) = \qty{468}{yr}$; और 0.5 हो, तो \qty{749}{yr}। अंतरापृष्ठ जितना कम अवरुद्ध है, एक महान भूकंप का सरकन जमा होने में उतना ही अधिक समय लेता है; यहाँ पाया गया औसत सरकन, \qtyproduct{400 x 150}{km} के आयत और $\mu = \qty{30}{GPa}$ के साथ,
स्वयं केवल एक अनुमान है, और वास्तविक सरकन लगभग समान नहीं था।
\end{solution}

\begin{solution}{exo:b1:earthquake-source:7}
(a) \omterm{def:b1:geological-maps:attitude}{स्ट्राइक} $\phi = 329$ है, \omterm{def:b1:geological-maps:attitude}{नति-दिशा} $a = 059$, \omterm{def:b1:geological-maps:attitude}{नति} $\delta =
\ang{8}$। तब
\[
  \begin{aligned}
    \mathbf s &= (\sin329, \cos329, 0) = (-0.5150, 0.8572, 0),\\
    \mathbf u &= (-\cos8\sin59, -\cos8\cos59, \sin8) = (-0.8488, -0.5100, 0.1392),\\
    \mathbf n &= (\sin8\sin59, \sin8\cos59, \cos8) = (0.1193, 0.0717, 0.9903)
  \end{aligned}
\]
और $r = \ang{110}$ के साथ,
\[
  \mathbf d = \cos110\,\mathbf s + \sin110\,\mathbf u = (-0.6215, -0.7724, 0.1308) .
\]
(b) \omterm{def:b1:earthquake-source:focal-mechanism}{सहायक तल} का अभिलंब $\mathbf d$ और उसका सरकन $\mathbf n$ है। ऊपर की ओर लिए गए अभिलंब के साथ, वह $\arccos 0.1308 = \ang{82.5}$ नत है,
$\arctan(-0.6215/-0.7724) = 218.8$ की ओर, इसलिए उसका \omterm{def:b1:geological-maps:attitude}{स्ट्राइक} $128.8$ है; उसके \omterm{def:b1:geological-maps:attitude}{स्ट्राइक} और \omterm{def:b1:geological-maps:attitude}{ऊपर-नति} सदिशों पर प्रक्षेपित सरकन
$\mathbf n$ \omterm{def:b1:faults-folds-fabrics:fault-elements}{रेक} \ang{87.2} देता है। पूर्णांकित: $129/83/87$, प्रकाशित हल का दूसरा तल।
(c) $\mathbf P = (\mathbf n - \mathbf d)/\sqrt2$ की दिगंश $221$ है और अवनति \ang{37};
$\mathbf T = (\mathbf n + \mathbf d)/\sqrt2$ की दिगंश $036$ है और अवनति
\ang{52} (प्रकाशित: $222/38$ और $36/52$)। (d) उथला तल, जो \ang{8} नत है, $059$ की ओर, \omterm{def:b1:active-margins:slab}{धँसते स्लैब} के साथ: \omterm{def:b1:faults-folds-fabrics:fault-elements}{रेक} \ang{110} वाला एक \omterm{def:b1:faults-folds-fabrics:fault-types}{उत्क्रम भ्रंश}, और \ang{45} से कम नत होने से, एक अधिक्षेप, सबडक्शन-क्षेत्र का प्लेट-अंतरापृष्ठ।
\end{solution}

\begin{solution}{exo:b1:earthquake-source:8}
(a) भ्रंश $135$ की ओर नत है, इसलिए $\mathbf n = (\sin135, \cos135, 0) =
(0.707, -0.707, 0)$; \omterm{def:b1:faults-folds-fabrics:fault-types}{दाएँ-पार्श्वीय} सरकन का $r = \ang{180}$ है, $\mathbf d = -\mathbf s
= (-0.707, -0.707, 0)$। उत्तर की ओर, $\mathbf x = (0, 1, 0)$:
$(\mathbf n\cdot\mathbf x)(\mathbf d\cdot\mathbf x) = (-0.707)(-0.707) > 0$, एक धक्का। पूर्व: $(0.707)(-0.707) < 0$, एक खिंचाव। दक्षिण: धक्का; पश्चिम: खिंचाव।
(b) $\mathbf P = (\mathbf n - \mathbf d)/\sqrt2 = (1, 0, 0)$: \omterm{def:b1:earthquake-source:focal-mechanism}{P-अक्ष} क्षैतिज है और
$090$--$270$ दिशा में है; $\mathbf T = (0, -1, 0)$ की दिशा
$000$--$180$ है। (c) $135$--$315$ \omterm{def:b1:geological-maps:attitude}{स्ट्राइक} करने वाले ऊर्ध्वाधर तल का अभिलंब
$\mathbf d$ है। सरकन $\mathbf n$ के साथ, उसके दक्षिण-पश्चिम की ओर का खंड, जिस ओर
$\mathbf d$ इंगित करता है, $135$ की ओर चलता है; उत्तर-पूर्व खंड से, दक्षिण-पश्चिम की ओर देखते हुए, वह बाईं ओर है: एक \omterm{def:b1:faults-folds-fabrics:fault-types}{बाएँ-पार्श्वीय} सरकन। \cref{prop:b1:earthquake-source:auxiliary} से यह सरकन हर दिशा में वही प्रथम गतियाँ देता है।
\end{solution}

\begin{solution}{exo:b1:earthquake-source:9}
(a) $V = \qty{40}{mm.yr^{-1}}$ और \qty{10}{km} पर $v = V/4$, इसलिए
\cref{prop:b1:earthquake-source:profile}(b) से $D = \qty{10}{km}$; सीधे,
$D = x/\tan(\pi v/V) = 10/\tan(\pi/4) = \qty{10}{km}$। (b) $v(30) =
(40/\pi)\arctan 3 = \qty{15.9}{mm.yr^{-1}}$। (c) $V/(\pi D) =
\num{0.040}/(\pi\times\num{10000}) = \num{1.27e-6}$ प्रति वर्ष।
\end{solution}

\begin{solution}{exo:b1:earthquake-source:10}
(a) $C q^k$ खानों के आघूर्ण अब भी $q = 10^{1.5 - b}$ के साथ हैं, परंतु अब
$q < 1$: खानों का आघूर्ण बड़े परिमाणों की ओर घटता है, और
निचला खाना, $k = 0$, $C/S = (1 - q)/(1 - q^{K+1}) > 1 - q$ का भाग वहन करता है, जो
\cref{prop:b1:earthquake-source:largest} का दर्पण-प्रतिबिंब है। गिने गए सबसे छोटे
भूकंप तब प्रभुत्व में होते, और जैसे-जैसे और छोटे भूकंप शामिल होते, कुल सीमा के बिना बढ़ता।
(b) $1 - 10^{-0.7} = 0.80$: कम से कम \qty{80}{\percent}। (c) $\log_{10}N = 8.20 - 1.01\times9 = -0.89$, इसलिए $N = 0.13$
प्रति वर्ष और पचास वर्षों में लगभग 6.4, जो प्रेक्षित दो से तीन गुना है। रेखा, जो परिमाण 6 और 7.5 के बीच आरोपित है, सबसे बड़े भूकंपों को अधिक आँकती है;
या तो नियम शीर्ष पर नीचे मुड़ता है, जैसा भ्रंशों का परिमित आकार सुझाता है, या पचास वर्ष दशकों में एक बार आने वाली घटनाओं के नमूने के लिए बहुत कम हैं। दो घटनाओं से दोनों संभव हैं।
\end{solution}

\begin{solution}{exo:b1:earthquake-source:11}
(a) 150 वर्षों में $x$ पर एक खंभा $150\,v(x) = 150\times(34/\pi)
\arctan(x/15)$ मिलीमीटर भ्रंश पर के खंभे के सापेक्ष चलता है: \qty{0.11}{m} पर
\qty{1}{km}, \qty{1.28}{m} \qty{15}{km} पर, \qty{2.08}{m} \qty{50}{km} पर,
दक्षिण की ओर, पश्चिम से देखे गए \omterm{def:b1:faults-folds-fabrics:fault-types}{दाएँ-पार्श्वीय} भ्रंश के लिए। (b) सरकन-घाटा है
$s = \qty{0.034}{m.yr^{-1}}\times\qty{150}{yr} = \qty{5.1}{m}$। भूकंप के दौरान खंभे $(s/\pi)\arctan(D/x)$ चलते हैं: \qty{2.44}{m} पर
\qty{1}{km}, \qty{1.28}{m} \qty{15}{km} पर, \qty{0.47}{m} \qty{50}{km} पर।
(c) योग $0.11 + 2.44$, $1.28 + 1.28$ और $2.08 + 0.47$ हैं, सब
\qty{2.55}{m} $= s/2$: भूकंप के बाद बाड़ पूर्वी ओर फिर सीधी है, $s/2$ विस्थापित, और पश्चिमी ओर, सममिति से, $-s/2$ से।
\end{solution}

\begin{solution}{exo:b1:earthquake-source:12}
(a) $\mathbf P = (0, 1, 0)$ और $\mathbf T = (1, 0, 0)$। युग्म
$\mathbf n = (\mathbf T + \mathbf P)/\sqrt2 = (0.707, 0.707, 0)$,
$\mathbf d = (\mathbf T - \mathbf P)/\sqrt2 = (0.707, -0.707, 0)$ संतुष्ट करता है
$\mathbf P = (\mathbf n - \mathbf d)/\sqrt2$ और $\mathbf T = (\mathbf n +
\mathbf d)/\sqrt2$। दोनों सदिश क्षैतिज हैं, इसलिए दोनों \omterm{def:b1:earthquake-source:focal-mechanism}{नोडल तल} ऊर्ध्वाधर हैं:
अभिलंब $\mathbf n$ वाला तल ($045$ की ओर) $135$--$315$ \omterm{def:b1:geological-maps:attitude}{स्ट्राइक} करता है,
अभिलंब $\mathbf d$ वाला तल ($135$ की ओर) $045$--$225$ \omterm{def:b1:geological-maps:attitude}{स्ट्राइक} करता है। (b) $315$ \omterm{def:b1:geological-maps:attitude}{स्ट्राइक} करने वाले तल पर, उत्तर-पूर्व की ओर का खंड $\mathbf d$ के अनुदिश चलता है,
$135$ की ओर: दूसरे खंड से देखने पर दाईं ओर, एक \omterm{def:b1:faults-folds-fabrics:fault-types}{दाएँ-पार्श्वीय} सरकन। $045$ \omterm{def:b1:geological-maps:attitude}{स्ट्राइक} करने वाले तल पर, दक्षिण-पूर्व की ओर का खंड $\mathbf n$ के अनुदिश
चलता है, $045$ की ओर: एक \omterm{def:b1:faults-folds-fabrics:fault-types}{बाएँ-पार्श्वीय} सरकन। (c) एक नया भ्रंश $\ang{45} - \phi/2 =
\ang{30}$ के साथ $\sigma_1$ बनाता है: $\sigma_1$ के $000$ दिशा में होने पर, भ्रंश $030$ (\omterm{def:b1:faults-folds-fabrics:fault-types}{बाएँ-पार्श्वीय}) और $330$ (\omterm{def:b1:faults-folds-fabrics:fault-types}{दाएँ-पार्श्वीय}) \omterm{def:b1:geological-maps:attitude}{स्ट्राइक} करेंगे, नोडल तलों से \ang{15} दूर।
\omterm{def:b1:earthquake-source:focal-mechanism}{P-अक्ष} ज्यामिति से नोडल तलों से \ang{45} पर तय है, चट्टान के घर्षण से स्वतंत्र, इसलिए वह $\sigma_1$ के साथ तभी मेल खाता है जब भ्रंश उससे \ang{45} पर सरका हो। (d) $040$ \omterm{def:b1:geological-maps:attitude}{स्ट्राइक} करने वाला भ्रंश $045$ \omterm{def:b1:geological-maps:attitude}{स्ट्राइक} करने वाले \omterm{def:b1:earthquake-source:focal-mechanism}{नोडल तल} से \ang{5} दूर है: वही तल संभवतः भ्रंश है, और उस पर सरकन \omterm{def:b1:faults-folds-fabrics:fault-types}{बाएँ-पार्श्वीय} है।
\end{solution}

\begin{solution}{pb:b1:earthquake-source:1}
\textbf{1.} अंतराल $1290 - 1050 = 240$, $1480 - 1290 = 190$ और
$1712 - 1480 = 232$ वर्ष हैं; उनका माध्य \qty{220.7}{yr} है।

\textbf{2.} औसत विस्थापन $(7.0 + 6.0 + 5.5 + 7.5)/4 = \qty{6.5}{m}$ है, और
$T = 6.5/0.030 = \qty{216.7}{yr}$, जो माध्य अंतराल के निकट है।

\textbf{3.} अंतिम तीन भूकंपों के विस्थापन, $6.0 + 5.5 + 7.5 =
\qty{19.0}{m}$, 1050 और 1712 के बीच 662 वर्षों में जमा हुए: $\qty{19.0}{m}
/\qty{662}{yr} = \qty{28.7}{mm.yr^{-1}}$, भू-आकलन के \qty{30}{mm.yr^{-1}} से \qty{4}{\percent} कम। दोनों मेल खाते हैं, अनिश्चितता के भीतर, तिथियों और विस्थापनों की: खंड अपने महान भूकंपों में अपने लगभग पूरे भार को मुक्त करता है।

\textbf{4.} $2025 - 1712 = 313$ वर्ष, और घाटा
$\qty{0.030}{m.yr^{-1}}\times\qty{313}{yr} = \qty{9.39}{m}$ है।

\textbf{5.} घाटा खंदक के हर विस्थापन से अधिक है, जिनमें सबसे बड़ा
\qty{7.5}{m} है, और 1712 से अब तक का समय सबसे लंबे अंतराल, 240 वर्ष, से अधिक है। इस अर्थ में खंड समय से पीछे है; इसका अर्थ यह नहीं कि भूकंप आसन्न है, केवल यह कि भ्रंश के पास अभिलेख में एक साथ मुक्त हुए सरकन से अधिक सरकन पहले से है।

\textbf{6.} अंतराल 190 से 240 वर्ष तक और विस्थापन \qtyrange{5.5}{7.5}{m} तक हैं: चक्र आवर्ती नहीं है, और अगले भूकंप पर भ्रंश की सामर्थ्य अज्ञात है। खंदक एक माध्य और एक प्रसार देती है, जिससे एक प्रायिकता, तिथि नहीं, आँकी जा सकती है।

\textbf{7.} $A = \qty{250}{km}\times\qty{15}{km} = \qty{3750}{km^2} =
\qty{3.75e9}{m^2}$।

\textbf{8.} $M_0 = \num{3e10}\times\num{3.75e9}\times\num{9.39} =
\qty{1.06e21}{N.m}$।

\textbf{9.} $M_0 = \num{3e10}\times\num{3.75e9}\times\num{7.5} =
\qty{8.44e20}{N.m}$, इसलिए $M_w = \tfrac23(20.926 - 9.1) = 7.88$।

\textbf{10.} अनुपात $9.39/7.5 = 1.25$ है, जो $\tfrac23\log_{10}1.25 =
0.07$ परिमाण है।

\textbf{11.} आधा क्षेत्रफल आघूर्ण को आधा करता है और परिमाण को
$\tfrac23\log_{10}2 = 0.20$ से घटाता है: $M_w = 7.75$।

\textbf{12.} $\mu A V = \num{3e10}\times\num{3.75e9}\times\num{0.030} =
\qty{3.38e18}{N.m}$ प्रति वर्ष।

\textbf{13.} $b = 1$। $N(5) = 10^{4.3 - 5} = 0.20$ प्रति वर्ष और $N(6) = 0.020$ प्रति
वर्ष।

\textbf{14.} $N(7) = 0.0020$ प्रति वर्ष: हर 500 वर्ष में एक; $N(8) = 0.00020$ प्रति
वर्ष: हर \num{5000} वर्ष में एक।

\textbf{15.} नियम परिमाण 7.95 या उससे अधिक का एक भूकंप हर
$1/10^{4.3 - 7.95} = \num{4500}$ वर्ष में देता है, खंदक हर 220 वर्ष में एक, यानी बीस गुना बार-बार। नियम क्षेत्र के सभी भ्रंशों के असंख्य छोटे भूकंपों पर आरोपित है; एक खंड के महान भूकंपों तक बहिर्वेशित किए जाने पर वह उन्हें बुरी तरह कम आँकता है। किसी भ्रंश पर महान भूकंपों की पुनरावृत्ति उसके अपने अभिलेख और उसके भार से ली जाती है, छोटों से नहीं।

\textbf{16.} $\num{5.5e17}/\num{3.38e18} = 0.16$: मध्यम भूकंप खंड द्वारा
संचित आघूर्ण का \qty{16}{\percent} मुक्त करते हैं।

\textbf{17.} \cref{prop:b1:earthquake-source:moment-steps} से,
$10^{1.5\times(7.95 - 6)} = 10^{2.92} \approx 840$ गुना अधिक।

\textbf{18.} $b = 1$ के साथ, भूकंपों का आघूर्ण सबसे बड़े से प्रभुत्व में होता है: मध्यम भूकंप जो भी करें, वे एक छोटा भाग, यहाँ \qty{16}{\percent}, मुक्त करते हैं, और शेष \qty{84}{\percent}, यानी सरकन-घाटा, अगले महान भूकंप की प्रतीक्षा करता है।

\textbf{19.} हर युग्म के लिए $v(-x) = -v(x)$, इसलिए परिच्छेद विषम है।
भ्रंश से दूर वेग $\pm V/2 = \pm\qty{15}{mm.yr^{-1}}$ की ओर प्रवृत्त होते हैं।

\textbf{20.} $D = x/\tan(\pi v/V) = 15/\tan(\pi\times7.5/30) = 15/\tan\ang{45} =
\qty{15}{km}$: \omterm{def:b1:earthquake-source:locking-depth}{अवरोधन-गहराई} पर वेग $V/4$ है।

\textbf{21.} \qty{5}{km} पर: $D = 5/\tan(\pi\times3.1/30) = 5/\tan\ang{18.6} =
\qty{14.9}{km}$; \qty{40}{km} पर: $D = 40/\tan\ang{69.6} = \qty{14.9}{km}$। वेग \qty{0.1}{mm.yr^{-1}} तक दिए गए हैं, और अकेला वह पूर्णांकन $D$ को लगभग \qty{0.1}{km} खिसका देता है; वास्तविक आँकड़ों के साथ हर स्टेशन का शोर अधिक होगा, और \omterm{def:b1:earthquake-source:locking-depth}{अवरोधन-गहराई} सभी स्टेशनों पर एक साथ निर्धारित की जाएगी।

\textbf{22.} \omterm{def:b1:earthquake-source:locking-depth}{अवरोधन-गहराई}, \qty{15}{km}, सबसे गहरे भूकंपों की गहराई के बराबर है। भ्रंश का अवरुद्ध भाग उसका भंगुर भाग है, \omterm{def:b1:stress-strain-rheology:strength-envelope}{भूकंप-जनक परत}, जहाँ घर्षण उसे भूकंपों के बीच रोकता है और उनमें विफल होता है; उसके नीचे, गरम चट्टान विसर्पित होती है (\cref{ch:b1:stress-strain-rheology}) और भ्रंश स्थिरता से सरकता है, यानी मॉडल में वेग $V$ पर, $D$ के नीचे सरकन।

\textbf{23.} $V/(\pi D) = \num{0.030}/(\pi\times\num{15000}) = \num{6.4e-7}$
प्रति वर्ष।

\textbf{24.} $M_0 = \qty{1.056e21}{N.m}$ के साथ, $\log_{10}M_0 = 21.024$ और
$M_w = \tfrac23(21.024 - 9.1) = 7.949$। 0.01 तक, वह भूकंप जो खंड का पूरा घाटा मुक्त करेगा, उसका \omterm{def:b1:earthquake-source:mw}{आघूर्ण परिमाण}
\textbf{7.95} है।
\end{solution}
